\documentclass[11pt,letterpaper]{article}

\usepackage[margin=1in]{geometry}
\usepackage{graphicx}
\usepackage{booktabs}
\usepackage{amsmath}
\usepackage{hyperref}
\usepackage{xcolor}
\usepackage{caption}
\usepackage{float}

\hypersetup{
  colorlinks=true,
  linkcolor=blue,
  citecolor=blue,
  urlcolor=blue
}

% --- Local macros (subset of PPG12 defs.sty) ---
\newcommand{\pp}{\mbox{$p$+$p$}}
\newcommand{\sqs}{\mbox{$\sqrt{s}$}}
\newcommand{\DR}{\ensuremath{\Delta R}}
\newcommand{\pT}{\ensuremath{p_{\text{T}}}}
\newcommand{\ET}{\ensuremath{E_{\text{T}}}}
\newcommand{\zvtx}{\ensuremath{z_{\mathrm{vtx}}}}
\newcommand{\Rcone}{\ensuremath{R_{\text{cone}}}}
\newcommand{\abseta}{\ensuremath{|\eta|}}
\newcommand{\isoETtruth}{\ensuremath{E_\text{T}^\text{iso,truth}}}

\graphicspath{{../../plotting/figures/deltaR_study/}}

\title{%
  \vspace{-1cm}
  {\large \textbf{sPHENIX Internal Note}}\\[0.5cm]
  {\LARGE \textbf{Study of the $\DR$ Truth-Matching Requirement\\
  in Single-Interaction MC}}\\[0.3cm]
  {\large Isolated Photon Cross-Section in \pp{} at $\sqs = 200$~GeV}
}
\author{sPHENIX Collaboration}
\date{\today}

\begin{document}

\maketitle

\begin{abstract}
This note documents a study of the $\DR < 0.1$ geometric truth-matching
requirement used in the reconstruction efficiency calculation for the sPHENIX
isolated photon cross-section measurement (PPG12). While a previous study of
double-interaction MC found a 65\% matching failure rate driven by vertex
displacement, the single-interaction samples were assumed to be unaffected.
We show that even in single-interaction MC, the $\DR$ cut rejects
$\sim$3.5\% of correctly track-ID-matched truth photons, with the rejection
driven entirely by vertex mismatch rather than cluster quality. 73\% of
rejected clusters have $\ET / \pT^{\text{truth}}$ in the range $[0.8, 1.2]$
and are well-reconstructed photons. The $\DR < 0.1$ cut is removed from the
production efficiency calculation; track ID matching alone provides the
definitive truth association.
\end{abstract}

\tableofcontents
\newpage

%% ============================================================
\section{Introduction}
\label{sec:intro}

The reconstruction efficiency for the sPHENIX isolated photon cross-section
measurement is computed from single-interaction PYTHIA + GEANT4 MC samples.
A truth photon is considered ``matched'' to a reconstructed cluster if two
criteria are satisfied:
\begin{enumerate}
  \item \textbf{Track ID match:} the GEANT track ID (\texttt{trkID}) of the
    truth particle matches that of a reconstructed cluster. This is the
    definitive truth association provided by the simulation framework.
  \item \textbf{Geometric match:} the angular separation between the truth
    photon and the matched cluster satisfies
    $\DR = \sqrt{(\Delta\eta)^2 + (\Delta\phi)^2} < 0.1$.
\end{enumerate}

The $\DR$ cut was introduced as a geometric sanity check on top of the
\texttt{trkID} association. A previous study of double-interaction MC
(documented in Ref.~\cite{di_note}) demonstrated that vertex displacement
causes a 65\% $\DR$-matching failure rate. In double-interaction events, the
reconstructed vertex is the average of two collision vertices, inducing large
$\Delta z$ shifts that propagate to $\Delta\eta$ via the EMCal geometry.

That study focused on double-interaction samples, and single-interaction MC was
assumed to be largely unaffected. This note quantifies the $\DR$ cut's impact
in single-interaction MC and evaluates whether it is needed at all.

\begin{thebibliography}{1}
\bibitem{di_note} sPHENIX Internal Note: Double-Interaction MC Efficiency
  Study, 2025.
\end{thebibliography}

%% ============================================================
\section{Method}
\label{sec:method}

The study uses the macro \texttt{study\_deltaR\_single.C}, processing the full
single-interaction MC samples used in the production efficiency calculation:
\texttt{photon5} (truth $\pT$ range 0--14~GeV), \texttt{photon10} (14--30~GeV),
and \texttt{photon20} (30+~GeV). Each sample is weighted by its production
cross-section to form a combined spectrum spanning the analysis range
$8 \leq \pT < 36$~GeV.

Truth photon selection follows the standard PPG12 criteria:
\begin{itemize}
  \item Particle ID: $\texttt{pid} = 22$
  \item Photon class: $\texttt{photonclass} < 3$ (direct photons)
  \item Truth isolation: $\isoETtruth < 4$~GeV (cone $\Rcone = 0.3$)
  \item Pseudorapidity: $\abseta < 0.7$
  \item Transverse momentum: $8 \leq \pT < 36$~GeV
\end{itemize}

For each truth photon, the closest \texttt{trkID}-matched reconstructed cluster
is identified, and $\DR$ is computed. The analysis scans over $\DR$ thresholds
of 0.05, 0.08, 0.10, 0.15, 0.20, 0.30, and no geometric cut (track ID only).
The \texttt{photon10\_double} sample is also processed for direct comparison
with the double-interaction results.

%% ============================================================
\section{Results}
\label{sec:results}

%% --------------------------------------------------
\subsection{\texorpdfstring{$\DR$}{Delta R} Failure Rates}
\label{sec:rates}

Table~\ref{tab:persample} summarizes the matching failure rates per sample and
for the combined (unweighted) dataset.

\begin{table}[htbp]
\centering
\caption{Truth-matching failure rates in single-interaction MC, per sample and
  combined. ``No trkID match'' indicates no reconstructed cluster carries the
  truth photon's track ID. ``Fail $\DR < 0.1$'' is the fraction of
  \texttt{trkID}-matched photons that fail the geometric cut.}
\label{tab:persample}
\begin{tabular}{lccc}
\toprule
\textbf{Sample} & \textbf{Truth photons} & \textbf{No trkID match} &
  \textbf{Fail $\DR < 0.1$ (of matched)} \\
\midrule
photon5  & 273{,}017   & 3.42\% & 2.70\% \\
photon10 & 2{,}380{,}500 & 2.37\% & 2.98\% \\
photon20 & 2{,}840{,}703 & 1.72\% & 4.01\% \\
\midrule
Combined & 5{,}494{,}220 & 2.09\% & 3.50\% \\
\bottomrule
\end{tabular}
\end{table}

The $\DR < 0.1$ cut rejects 2.7--4.0\% of truth photons that already have a
valid \texttt{trkID} match, with the rate increasing at higher $\pT$ where
forward-$\eta$ clusters are more sensitive to vertex shifts.

Table~\ref{tab:singledouble} compares the single-interaction and
double-interaction results for the \texttt{photon10} sample, illustrating the
continuity between the two regimes.

\begin{table}[htbp]
\centering
\caption{Comparison of \texttt{photon10} single-interaction and
  double-interaction MC. The $\Delta z$ RMS is the event-level distribution
  of $|z_{\text{vtx}}^{\text{reco}} - z_{\text{vtx}}^{\text{truth}}|$.}
\label{tab:singledouble}
\begin{tabular}{lccc}
\toprule
\textbf{Sample} & \textbf{$\Delta z$ RMS} & \textbf{No trkID} &
  \textbf{Fail $\DR < 0.1$ (of matched)} \\
\midrule
photon10 single & 9.8~cm & 2.37\% & 3.06\% \\
photon10 double & 46.2~cm & 7.92\% & 65.0\% \\
\bottomrule
\end{tabular}
\end{table}

The $\DR$ failure rate scales with vertex displacement: the double-interaction
sample's much larger $\Delta z$ RMS drives a proportionally larger failure
rate, consistent with the geometric mechanism established in the
double-interaction study.

\begin{figure}[htbp]
\centering
\includegraphics[width=0.85\textwidth]{deltaR_distribution.pdf}
\caption{Distribution of $\DR$ between truth photons and their
  \texttt{trkID}-matched reconstructed clusters in single-interaction MC.
  The vertical line indicates the $\DR = 0.1$ threshold. A tail of events
  extends beyond the cut, corresponding to vertex-displaced clusters.}
\label{fig:dr_dist}
\end{figure}

%% --------------------------------------------------
\subsection{Threshold Scan}
\label{sec:scan}

Table~\ref{tab:scan} shows the $\DR$ failure rate as a function of the
matching threshold for both single and double-interaction \texttt{photon10}
samples.

\begin{table}[htbp]
\centering
\caption{Fraction of \texttt{trkID}-matched truth photons failing the $\DR$
  cut as a function of threshold, for \texttt{photon10} single and double
  interaction MC.}
\label{tab:scan}
\begin{tabular}{lcc}
\toprule
\textbf{Threshold} & \textbf{Single fail rate} & \textbf{Double fail rate} \\
\midrule
$\DR \geq 0.05$ & 5.03\% & 72.6\% \\
$\DR \geq 0.08$ & 3.60\% & 68.0\% \\
$\DR \geq 0.10$ & 3.06\% & 65.0\% \\
$\DR \geq 0.15$ & 2.27\% & 57.8\% \\
$\DR \geq 0.20$ & 1.83\% & 50.9\% \\
$\DR \geq 0.30$ & 1.27\% & 37.9\% \\
\bottomrule
\end{tabular}
\end{table}

Even relaxing the cut to $\DR < 0.30$ still rejects 1.3\% of single-interaction
photons, indicating that the issue is not a matter of tuning the threshold
but reflects the fundamental limitation of a geometric cut applied on top of a
vertex-dependent coordinate system.

\begin{figure}[htbp]
\centering
\includegraphics[width=0.85\textwidth]{efficiency_threshold_scan.pdf}
\caption{Matching efficiency as a function of $\DR$ threshold for
  \texttt{photon10} single-interaction (blue) and double-interaction (red) MC.
  The nominal threshold of 0.1 is marked. In single-interaction MC, removing
  the cut entirely recovers $\sim$3\% efficiency.}
\label{fig:scan}
\end{figure}

%% --------------------------------------------------
\subsection{Vertex Mismatch as the Driving Mechanism}
\label{sec:vertex}

The connection between vertex mismatch and $\DR$ failure is examined by
computing the failure fraction as a function of
$|\Delta z| = |z_{\text{vtx}}^{\text{reco}} - z_{\text{vtx}}^{\text{truth}}|$.
Table~\ref{tab:dz} shows a sharp step function.

\begin{table}[htbp]
\centering
\caption{$\DR < 0.1$ failure fraction as a function of $|\Delta z|$ in
  single-interaction MC. The transition occurs at the predicted critical
  threshold $|\Delta z|_{\text{crit}} = 9.35$~cm.}
\label{tab:dz}
\begin{tabular}{lc}
\toprule
\textbf{$|\Delta z|$ range (cm)} & \textbf{Failure fraction} \\
\midrule
$< 9$   & $\sim$0\%  \\
9--10  & 0.9\%   \\
10--11 & 20\%    \\
11--12 & 62\%    \\
12--13 & 88\%    \\
$\geq 15$ & 100\%   \\
\bottomrule
\end{tabular}
\end{table}

The critical threshold is predicted by the EMCal barrel geometry:
\begin{equation}
  |\Delta z|_{\text{crit}} = \DR_{\text{cut}} \times R_{\text{EMCal}}
    \times \cosh\eta
    = 0.1 \times 93.5~\text{cm}
    = 9.35~\text{cm} \quad (\text{at } \eta = 0)
  \label{eq:dzcrit}
\end{equation}
The observed step at $|\Delta z| \approx 9$--10~cm matches this prediction
precisely, confirming that vertex mismatch is the sole driver of the failure.

In the single-interaction MC, the cross-section-weighted $\Delta z$ RMS is
2.88~cm, but the vertex distribution has a long tail: 2.12\% of truth photons
have $|\Delta z| > 9.35$~cm. Virtually all of these fail the $\DR < 0.1$ cut.

\begin{figure}[htbp]
\centering
\includegraphics[width=0.85\textwidth]{fail_fraction_vs_dz.pdf}
\caption{$\DR < 0.1$ failure fraction as a function of $|\Delta z|$ in
  single-interaction MC. The sharp step at $\sim$9.35~cm matches the
  geometric prediction of Eq.~(\ref{eq:dzcrit}).}
\label{fig:fail_dz}
\end{figure}

\begin{figure}[htbp]
\centering
\includegraphics[width=0.85\textwidth]{deta_vs_dz.pdf}
\caption{Correlation between $\Delta\eta$ (truth$-$reco) and vertex shift
  $\Delta z$. The linear band follows $\Delta\eta \approx
  -\Delta z / (R_{\text{EMCal}} \cdot \cosh\eta)$, confirming the geometric
  origin of the $\DR$ failure.}
\label{fig:deta_dz}
\end{figure}

\begin{figure}[htbp]
\centering
\includegraphics[width=0.85\textwidth]{deltaR_vs_dz.pdf}
\caption{Two-dimensional distribution of $\DR$ vs.\ $|\Delta z|$ in
  single-interaction MC. Events above the horizontal line ($\DR = 0.1$) are
  rejected by the matching cut; they cluster at $|\Delta z| > 9$~cm.}
\label{fig:dr_dz}
\end{figure}

%% --------------------------------------------------
\subsection{Energy Quality of Rejected Clusters}
\label{sec:quality}

If the $\DR$ cut were removing genuinely mis-reconstructed clusters, we would
expect degraded energy response. Figure~\ref{fig:etratio} shows the cluster
$\ET / \pT^{\text{truth}}$ distribution separately for clusters that pass and
fail the $\DR < 0.1$ cut.

\begin{table}[htbp]
\centering
\caption{Cluster $\ET / \pT^{\text{truth}}$ distribution for
  \texttt{trkID}-matched clusters passing and failing the $\DR < 0.1$ cut.}
\label{tab:quality}
\begin{tabular}{lcc}
\toprule
\textbf{Category} & \textbf{Mean $\ET / \pT^{\text{truth}}$} &
  \textbf{Fraction in $[0.8, 1.2]$} \\
\midrule
Pass ($\DR < 0.1$) & 0.945 & 88\% \\
Fail ($\DR \geq 0.1$) & 0.899 & 73\% \\
\bottomrule
\end{tabular}
\end{table}

The majority of rejected clusters are well-reconstructed photons. Their
calorimeter energy deposit is correct; only the angular position is shifted
by the vertex mismatch. The modest reduction in mean $\ET / \pT^{\text{truth}}$
for failing clusters (0.899 vs.\ 0.945) is itself a consequence of the
$\eta$-shift: the $\ET = E / \cosh\eta$ calculation uses the displaced vertex,
producing a slightly different $\ET$ than the truth-level $\pT$.

\begin{figure}[htbp]
\centering
\includegraphics[width=0.85\textwidth]{ET_ratio_pass_fail.pdf}
\caption{Distribution of cluster $\ET / \pT^{\text{truth}}$ for
  \texttt{trkID}-matched clusters that pass (blue) and fail (red) the
  $\DR < 0.1$ cut. 73\% of failing clusters have $\ET / \pT^{\text{truth}}
  \in [0.8, 1.2]$, indicating they are correctly reconstructed photons.}
\label{fig:etratio}
\end{figure}

%% --------------------------------------------------
\subsection{Extension to Double-Interaction MC: \texorpdfstring{$\Delta\eta$ vs.\ $\Delta z$}{Delta-eta vs.\ dz}}
\label{sec:double_deta_dz}

The vertex mismatch mechanism identified in single-interaction MC
(Section~\ref{sec:vertex}) is dramatically amplified in
double-interaction events. Figure~\ref{fig:deta_dz_comp} compares the
$\Delta\eta$ vs.\ $\Delta z$ correlation for \texttt{photon10}
single-interaction (left) and \texttt{photon10\_double} full GEANT
double-interaction MC (right), using 2M events per sample with the
standard truth photon selection ($\texttt{pid}=22$, $\texttt{photonclass}<3$,
$\isoETtruth < 4$~GeV, $\abseta < 0.7$, $8 \leq \pT < 36$~GeV).

In double-interaction events, the MBD-reconstructed vertex is
approximately the average of two collision vertices:
\begin{equation}
  z_{\text{vtx}}^{\text{reco}} \approx \frac{z_1 + z_2}{2}, \qquad
  \Delta z \approx \frac{z_2 - z_1}{2}
  \label{eq:double_dz}
\end{equation}
where $z_1$ is the primary (hard photon) collision vertex and $z_2$ is the
pileup collision vertex. This broadens the event-level $\Delta z$
distribution from an RMS of 16.9~cm (single, before vertex cut) to
48.1~cm (double), and correspondingly the cluster $\Delta\eta$
distribution from RMS 0.073 to 0.365.

Both distributions follow the geometric prediction
$\Delta\eta \approx -\Delta z / R_{\text{CEMC}}$ (red dashed line),
confirming that the $\eta$-shift mechanism is identical in both
regimes --- only the magnitude of the vertex displacement differs.

\begin{table}[htbp]
\centering
\caption{Summary of vertex displacement and matching statistics for
  \texttt{photon10} single and double-interaction MC (2M events each,
  trkID-matched truth photons within $|z_{\text{vtx}}^{\text{reco}}| < 60$~cm).}
\label{tab:double_summary}
\begin{tabular}{lcc}
\toprule
\textbf{Metric} & \textbf{Single} & \textbf{Double} \\
\midrule
Event-level $\Delta z$ RMS (all events) & 16.9~cm & 48.1~cm \\
Cluster $\Delta\eta$ RMS (matched) & 0.073 & 0.365 \\
No trkID match & 2.33\% & 7.95\% \\
Fail $\DR < 0.1$ (of matched) & 3.04\% & 65.08\% \\
Events with $|\Delta z| > 9.35$~cm & 36.0\% & 74.2\% \\
\bottomrule
\end{tabular}
\end{table}

The practical consequence: 65.1\% of trkID-matched clusters in
double-interaction MC fail the $\DR < 0.1$ cut (vs.\ 3.0\% in single),
and 74.2\% of events exceed the critical $|\Delta z| = 9.35$~cm
threshold. This provides further quantitative support for removing the
geometric matching requirement.

\begin{figure}[htbp]
\centering
\includegraphics[width=\textwidth]{deta_vs_dz_comparison.pdf}
\caption{Correlation between $\Delta\eta$ (cluster $-$ truth) and
  vertex displacement $\Delta z = z_{\text{vtx}}^{\text{reco}} -
  z_{\text{vtx}}^{\text{truth}}$ for \texttt{photon10}
  single-interaction MC (left) and full GEANT double-interaction MC
  (right). The red dashed line shows the geometric prediction
  $\Delta\eta = -\Delta z / R_{\text{CEMC}}$ at $\eta = 0$. Note the
  different axis scales: the double-interaction $\Delta z$ extends to
  $\pm 120$~cm (vs.\ $\pm 30$~cm for single) and $\Delta\eta$ to
  $\pm 1.2$ (vs.\ $\pm 0.15$).}
\label{fig:deta_dz_comp}
\end{figure}

%% ============================================================
\section{\texorpdfstring{Physics Argument for Removing the $\DR$ Cut}{Physics Argument for Removing the Delta R Cut}}
\label{sec:argument}

Six arguments support removing the $\DR < 0.1$ requirement:

\begin{enumerate}
  \item \textbf{Track ID is ground truth.} The GEANT \texttt{trkID} provides
    a definitive link between the truth particle and its energy deposit in the
    calorimeter. The $\DR$ cut adds no additional truth information; it tests
    a derived geometric quantity that depends on the vertex reconstruction.

  \item \textbf{The cut tests vertex resolution, not calorimeter performance.}
    The failure rate correlates with $|\Delta z|$ (Table~\ref{tab:dz}),
    not with any measure of cluster quality. The sharp step at
    $|\Delta z| = 9.35$~cm is a geometric artifact of the EMCal radius and the
    $\DR$ threshold.

  \item \textbf{Rejected clusters are good photons.} 73\% of failing clusters
    have $\ET / \pT^{\text{truth}} \in [0.8, 1.2]$
    (Table~\ref{tab:quality}). The calorimeter correctly measured the photon
    energy; only the vertex-dependent angular coordinate is shifted.

  \item \textbf{The response matrix should capture vertex-induced shifts.}
    Kinematic migration from vertex mismatch (the slight $\ET$ shift visible
    in the failing clusters) is precisely the kind of detector effect that the
    unfolding response matrix is designed to correct. Excluding these events
    from the efficiency numerator biases the response matrix.

  \item \textbf{The $\sim$3\% bias is vertex-dependent.} Because the failure
    rate depends on $|\Delta z|$, it is correlated with beam conditions
    (vertex spread, crossing angle). An efficiency correction that includes
    this bias would need to be parameterized in $\Delta z$, adding unnecessary
    complexity.

  \item \textbf{Consistency between single and double MC.} Removing the $\DR$
    cut makes the single-interaction and double-interaction efficiency
    calculations more directly comparable, simplifying the pileup systematic
    uncertainty evaluation.
\end{enumerate}

We note that the production code contains a commented-out alternative cut,
\texttt{(cluster\_ET / particle\_Pt) > 0.8}, which directly tests
calorimeter energy response rather than angular geometry. If any quality gate
beyond \texttt{trkID} is desired, an energy-response cut is more physically
motivated than a geometric $\DR$ cut, as it removes genuinely
mis-reconstructed clusters (merges, splits, lost energy) without penalizing
vertex displacement.

%% ============================================================
\section{Conclusion}
\label{sec:conclusion}

The $\DR < 0.1$ geometric truth-matching cut is removed from the production
reconstruction efficiency calculation. The key findings are:

\begin{enumerate}
  \item In single-interaction MC, the $\DR < 0.1$ cut rejects $\sim$3.5\% of
    \texttt{trkID}-matched truth photons. This is not negligible at the
    precision targets of the measurement.

  \item The rejection is driven entirely by vertex mismatch
    ($|\Delta z| > 9.35$~cm), not by cluster quality. A sharp step function
    in failure rate vs.\ $|\Delta z|$ matches the geometric prediction from
    the EMCal barrel radius.

  \item 73\% of rejected clusters are well-reconstructed photons with
    $\ET / \pT^{\text{truth}} \in [0.8, 1.2]$.

  \item Track ID matching alone is sufficient and definitive for truth
    association. The $\DR$ cut adds no truth information and introduces a
    vertex-dependent efficiency bias.

  \item The $\sim$3\% difference in efficiency can be assigned as a systematic
    uncertainty if desired, though the physics argument supports treating the
    track-ID-only matching as the correct baseline.

  \item The $\Delta\eta$ vs.\ $\Delta z$ correlation in full GEANT
    double-interaction MC (Section~\ref{sec:double_deta_dz}) confirms
    that the same geometric mechanism operates at larger vertex
    displacements ($\Delta z$ RMS 48.1~cm vs.\ 16.9~cm), with 65\% of
    matched clusters exceeding the $\DR = 0.1$ threshold.
\end{enumerate}

\end{document}
